
Unless flagged otherwise, every citation below offers adjacent, analogous, or precedent evidence for a construct \tphe{} uses --- never evidence that \tphe{} itself, or any of its five rules, works; this single sentence replaces a per-subsection restatement of that caveat.

\subsection{Burden}
The population-level case for treating the marriage transition as a safeguard point rests on burden-of-disease and surveillance data, not on any evaluation of premarital education. \S1.3 established the burden case (GBD 2023 \cite{ref1}; WHO/LSHTM \cite{ref2}: 27\% lifetime prevalence, UI 23--31\%, ever-partnered women 15--49, substantial exposure already in the 15--19 and 19--24 bands). A companion burden-of-proof analysis found a moderate association with major depressive disorder and maternal abortion/miscarriage (its largest), and only weak associations with HIV, anxiety and self-harm \cite{ref21}. A Lancet Psychiatry Commission situates these estimates within a wider call for mental-health-service investment \cite{ref22}; a Lancet Public Health review identified only 30 unique interventions from 16 countries addressing IPV and/or violence against children \cite{ref23} --- sparse relative to the burden. These are surveillance/modelling estimates, not evidence any intervention changes them.

\subsection{Primary prevention frameworks, what works, and recognition}
The strongest experimental evidence that structured relationship curricula can reduce IPV comes from community-mobilization and couples-curriculum trials outside the premarital-education literature. SASA! was associated with lower social acceptance of IPV (aRR 0.54) and a physical-IPV effect whose CI crossed 1 (aRR 0.48) \cite{ref24}, with norms change --- not knowledge or attendance --- explaining 70--95\% of the mediated effect \cite{ref25}. Indashyikirwa in Rwanda found a large reduction in women's reports of IPV (aRR 0.44) and men's reports of perpetration (aRR 0.54) \cite{ref26}, while a faith-based trial in Nigeria reported a 61\% reduction in odds of any reported IPV \cite{ref27} --- the clearest faith-setting positive result in this review. Set against these are equally rigorous null results: a scaled adaptation in rural Rwanda found no effect on any primary outcome \cite{ref28}, and a faith-leader caregiving trial in Kenya found no effect on its declared primary outcome while showing secondary IPV effects \cite{ref29} --- the outcome-substitution pattern \tphe{}'s failure criterion is designed to flag. Only 6/178 (3\%) reviews updating the WHO RESPECT framework targeted relationship-skills strengthening specifically \cite{ref30}, and a scoping review of implementation science found fidelity, acceptability and feasibility the most-studied outcomes, with scale, sustainability and survivor-centred practice minimally addressed \cite{ref31} --- a field-level gap consistent with treating \tphe{} as a programme theory rather than advancing an effectiveness claim at this stage.

Psychological inoculation --- drawn from misinformation research, not interpersonal coercion, and used here only as an analogy --- is the closest experimental literature to Layer 0 (a proposed extension)'s ``awareness as protection'' claim. The largest field experiment (7 preregistered studies, 6 RCTs, $n=6{,}464$, plus a YouTube field study $n=22{,}632$) found inoculation videos improved recognition of manipulation techniques and sharing-decision quality \cite{ref182} (effect sizes not reported in the abstract), extended to vaccine misinformation \cite{ref183}, election misinformation across 12 EU nations \cite{ref184}, a gamified format \cite{ref185}, real-world cross-protection against untrained techniques \cite{ref186}, and methodological work isolating technique from testing effects \cite{ref187} --- none of it about intimate-partner coercion. The strongest cautionary evidence comes from a 1-year cluster RCT of Coaching Boys Into Men (16 high schools, $n=1{,}513$ male athletes): perpetration fell, but recognition of abusive behaviours did not change (95\% CI $-$0.15 to 0.09) and bystander intentions were unchanged \cite{ref190} --- recognition and behaviour dissociating exactly as the failure criterion is built to catch. A Cochrane protocol for dating-violence prevention has not yet reported outcomes \cite{ref188}; a pilot questions whether attitudes proxy behaviour \cite{ref189}; further design and implementation papers report no recognition outcomes yet \cite{ref191,ref192}; a non-randomised bystander study found effects on risk recognition and perceived self-efficacy \cite{ref193}. Qualitatively, women describe difficulty naming gaslighting despite describing its effects \cite{ref194}; being silenced through negative responses to help-seeking, and loneliness, form part of the pathway to later intimate-partner violence and abuse in young adults \cite{ref195}; disclosure barriers are documented among sexual and gender-minority survivors \cite{ref196}; two studies of South and Southeast Asian women describe culturally specific understandings and barriers relevant to naming abuse \cite{ref200,ref201} (measurement tools not yet applied here: an internalised-stigma scale \cite{ref197}, Thai/general health-literacy instruments \cite{ref198,ref199}). This literature makes Layer 0's premise plausible and testable, not established; recognition and help-seeking are shown here to be separable outcomes.

\subsection{Preconception medicine and premarital screening as a proposed analogue}
Preconception medicine treats the period before conception/marriage as a distinct clinical window: the FIGO Preconception Checklist frames structured health-assessment elements for it \cite{ref11}; couple-based preconception-care programmes have been reviewed across 14 studies in 10 countries, mostly high-income, heterogeneous designs \cite{ref32}. Premarital genetic screening is the closest direct precedent for a premarital contact point at population scale (roughly 30--40\% of the Thai population carries a thalassaemia gene \cite{ref117}, motivating national couple-screening programmes): premarital sickle-cell-trait screening uptake across Africa pools under half of eligible couples (47.82\%, 95\% CI 46.53--49.11), ``suboptimal'' \cite{ref12}; premarital thalassaemia education produces consistent knowledge gains ($p<0.001$) but inconsistent screening-uptake behaviour \cite{ref13} --- the same knowledge/behaviour gap recurring below.

\subsection{Premarital/relationship education outcomes and the knowledge-to-well-being gap}
Across three cumulative meta-analyses spanning 1997--2022, the MRE evidence base has consistently measured relationship quality and communication skill \cite{ref6,ref33}, and more recently mental health and coparenting \cite{ref7}, with effects real but small ($d\approx 0.03$--$0.45$ \cite{ref6,ref7}) and authors flagging thin racial/ethnic diversity \cite{ref6} and rarely measured policy-relevant outcomes \cite{ref7}; a third meta-analysis in this lineage examined programmatic moderators such as dosage and communication emphasis \cite{ref33}. Individual RCTs complicate rather than resolve this: an 8-year follow-up found no divorce benefit and a harmful-moderation pattern for higher-conflict couples \cite{ref9}; a 3-year RCT found couples with acute baseline risk benefited less \cite{ref10}; a later meta-analysis of relationship aggression (31 studies, $n=25{,}527$) found $d\approx 0.11$ overall, not statistically convincing when restricted to RCTs ($d\approx 0.05$) \cite{ref8}. Newer, smaller trials sharpen the specific knowledge-to-well-being gap this paper labels its ``middle of the bridge'': a quasi-experimental study among engaged couples in Babol, Iran, found booklet- versus video-based premarital education improved reproductive-health-literacy access, evaluation and decision-making dimensions, but not reading-and-understanding \cite{ref34}; a 2026 study of reproductive/sexual-health education in premarital counselling classes in Rasht, Iran, found knowledge and attitudes improved at 3 months but reported no clinical outcomes \cite{ref35}; and a meta-analysis of online relationship-education programmes (study populations not characterised in the abstract) found they improved satisfaction, communication, confidence and individual-level anxiety/depressive symptoms \cite{ref36}. At the far end of this pathway, a large meta-analysis (126 articles, $>72{,}000$ participants) found marital quality associated with physical health, lower mortality risk and lower cardiovascular reactivity during conflict \cite{ref37} --- a well-being pathway exists in principle, not evidence premarital education produces the marital quality that pathway depends on. None of this literature has measured informed choice, coercion recognition, or safe help-seeking. Two frequently cited works in this literature --- Fawcett et al. (2010, \emph{Family Relations}) and Carroll \& Doherty (2003) --- are not indexed in PubMed; Fawcett et al. (2010) is cited once, flagged as not PubMed-indexed, in the novelty comparison (Table~\ref{tab:novelty}), and Carroll \& Doherty (2003) is not cited under this review's PubMed-only scope.

\subsection{Screening, first-line response, referral, and technology-facilitated routes}
The clinical screening literature offers the clearest precedent for the distinction \tphe{}'s failure criterion is built on. A Cochrane review of 13 RCTs found screening increased identification (OR 2.95) but no evidence of an effect on referral (OR 2.24, 2 studies only) or on re-exposure at 3--18 months \cite{ref38}. The 2025 USPSTF evidence report reviewed 35 studies ($n=18{,}358$): 3 RCTs ($n=3{,}759$) comparing screening with no screening found no significant reduction in IPV or other outcomes over 3--18 months, instrument sensitivity ranged 26--87\% (specificity 80--97\%), and of 13 intervention RCTs only perinatal home-visiting and multi-risk behavioural counselling showed benefit \cite{ref39}; the companion 2025 recommendation statement retained a moderate-net-benefit judgement for reproductive-age women \cite{ref40}. Provider training may improve knowledge, attitudes and self-perceived readiness, with weak and inconsistent evidence for actual response behaviour or survivor outcomes \cite{ref41}. Where first-line response protocols have been trialled directly (WHO LIVES/ARCHES in Nigeria), results are more encouraging but bounded: reduced relative risk of IPV exposure and reproductive coercion among women not already experiencing violence, but no effect on contraceptive use \cite{ref42}, with WHO's own guidance conditioning any inquiry about violence on a functioning referral system already in place \cite{ref43,ref44}. Premarital education itself may function as a gateway rather than a substitute for later help-seeking --- one US study found it associated with an increased likelihood of later couples counselling, with `not knowing where to go' reported alongside cost as a principal barrier among low-income couples who did not seek help \cite{ref106,ref107}. Delivery agents themselves may complicate a referral-first safeguard: in a survey of Christian clergy and congregation members, clergy reported being comfortable making referrals for professional counselling, whereas most congregation members preferred counselling with their own pastor and respondents wanted more training on domestic violence --- a live tension, not settled support, for a policy that assumes faith-based delivery agents will refer rather than counsel \cite{ref118}.

Evidence on Layer 2 (a proposed extension)'s moderated peer community is genuinely two-sided. On support: a review of youth mental-health peer forums found only 2 of 3 RCTs significant, the other four null \cite{ref139}; realist syntheses across mental-health and cancer settings agree outcomes depend on moderation, trust and anonymity/privacy control, noting ``not all online support groups are helpful, and in some cases they may increase distress'' \cite{ref140,ref141,ref146}. A scoping review of moderator practice found only 5 of 397 screened articles eligible, only one using a qualified health professional as moderator \cite{ref142}; qualitative work describes the role as under-supported and under-specified \cite{ref143}, while other qualitative work documents moderators' facilitative strategies and the need for better training and supervision \cite{ref144}. Weakly moderated platforms specifically amplified harmful content relative to strongly moderated ones \cite{ref145}. A Chinese IPV-specific peer-support study reports mostly supportive, rarely blaming, content \cite{ref147}, but existing online-IPV tools focus on leaving a relationship with little post-departure attention \cite{ref148}, and digital-IPV-tool safety research remains ``virtually non-existent'' \cite{ref137}. On harm: technology-facilitated abuse maps onto every domain of the Power and Control Wheel \cite{ref149}, and more broadly spans digital dating abuse, cyberstalking, technology-facilitated sexual assault, image-based sexual abuse, and online sexual harassment \cite{ref150}; immigrant/refugee women face additional, unique barriers to escaping and reporting it \cite{ref151}; it includes systematically characterised cyberstalking \cite{ref152}; and record-keeping features in coparenting apps have been shown to be sites of both coercive control and survivor resistance \cite{ref153}. Anonymity increases self-disclosure, while other features of online communication such as asynchronicity, multiple audiences, and audience feedback favor impression management over genuine self-presentation \cite{ref154}; Muslim youth's social media use carries both cultural/religious connection and discrimination \cite{ref157}; young people's reported benefits of online help-seeking come with literacy and trust barriers \cite{ref156}. This symmetry is why Layer 2 needs a named detection mechanism, not only a withdrawal rule.

\subsection{Couple-based work, coercive control, and assessor independence}
The evidence base most directly relevant to \tphe{}'s stop rule (R4) is small and explicitly cautious. The only systematic review of couples therapy as an IPV treatment, after screening 1{,}733 citations, found only 6 studies met inclusion; pooled data suggested conjoint treatment ``viable in select situations,'' with professionals cautious about violent-retaliation risk \cite{ref45}. Clinical literature converges on screening, safety modification, ongoing reassessment, and distinguishing situational couple violence from coercive controlling violence as prerequisites for any conjoint work \cite{ref46,ref47}. Survivor-voiced evidence complicates any default toward joint response --- only 9 of 17 women in one qualitative sample rated couples counselling a good idea \cite{ref48} --- and measurement is contested: coercive control is often reported as reciprocal \cite{ref49}, standard instruments do not measure perpetration and victimisation symmetrically \cite{ref50}, and even validated disclosure instruments substantially undercount coercive control through item-interpretation confusion and fear \cite{ref51}.

No study tests \tphe{}'s no-stake assessor (a proposed extension); the closest matched evidence is regulated independence in other high-stakes consent settings, offered here as analogy, not domain evidence. Transplant medicine mandates an Independent Living Donor Advocate separated from the recipient's care team \cite{ref158,ref159}, extended by proposal to two temporally separated consent processes for domino donors \cite{ref160}; facial plastic surgery literature names the same conflict-of-interest problem without a mandated remedy \cite{ref161}, and psychotherapy literature names the same dual-relationship problem, noting only meager research and little guidance on nonsexual dual relationships \cite{ref162}, with common occurrence in small or rural settings noted elsewhere \cite{ref163}. Structured IPV risk-assessment tools show independence helps but does not guarantee accuracy: a meta-analytic review found ODARA best-performing among five compared (weighted AUC=.666), yet only 9 of 20 (45\%) predictive-validity estimates used the tool exactly as intended, with a separate study validating an authorised German translation of the ODARA instrument \cite{ref164,ref168}; the Danger Assessment and its 5-item version retain real predictive validity for severe/lethal re-assault (AUC .68--.78) \cite{ref166}; the Taiwan-adapted version shows similarly strong discriminant and predictive validity, particularly for high dangerousness, though no specific AUC value is reported in the abstract \cite{ref165}; and a version updated to account for multiple/near-loss-of-consciousness strangulation predicts near-fatal re-assault significantly better than the original Danger Assessment, though no AUC value is reported in the abstract \cite{ref167}; but multiagency structured-judgment review found it often unclear which recorded factors drove a risk categorisation \cite{ref169}; a police-administered SARA field study found intervention intensity could reduce recidivism in high-risk cases while increasing it in low-risk cases \cite{ref170}; structured-judgment ratings, not only actuarial scores, differed by ethnicity between Indigenous and non-Indigenous participants in one study \cite{ref171}. Check-in and brief-contact models show the same mixed pattern the failure criterion is designed to catch: quality-improvement warm-handoff work doubled successful staff-to-community-health-worker referrals \cite{ref172}, yet an emergency-department screening programme found only 24.5\% of screened patients recalled receiving a resource guide despite 58.5\% being screened \cite{ref173}. An SMS brief-contact trial ($n=804$) found a 22\% relative reduction in self-harm repetition at 12 and 24 months but no difference in time to first repeat event \cite{ref174}, and a caring-contacts trial found no significant difference in loneliness at 6 months by treatment arm, concluding an added phone call increased complexity without improving effectiveness \cite{ref176} (protocol for the related comparative-effectiveness trial: \cite{ref175}). Implementation-fidelity work in an adjacent prevention programme shows that measured delivery fidelity and dose do not by themselves guarantee outcome change \cite{ref179}. This literature makes assessor independence necessary but not sufficient, and supports treating the day itself as presence of the space rather than as the outcome --- consistent with the failure criterion, not a validation of it.

\subsection{Decision quality, autonomy measurement, and digital self-assessment}
Decision quality --- shared decision-making and informed, values-congruent choice --- is this paper's proposed nearest measurement family for \tphe{}'s decision-control construct (not a validated instrument for it). The ``three-talk'' model structures shared decision-making into team, option and decision talk \cite{ref15}; a Cochrane review of 209 decision-aid studies ($>107{,}000$ participants) found consistent improvement in knowledge, risk-perception accuracy and values-congruent choice, and reduced decisional conflict \cite{ref14}. The MRC complex-intervention framework \cite{ref52}, its ADAPT context-adaptation guidance \cite{ref53}, and its process-evaluation guidance \cite{ref54} supply the methodological scaffolding for describing \tphe{} as a programme theory awaiting adaptation and process evaluation, not efficacy testing. The reproductive-autonomy and coercion instruments reviewed above \cite{ref16,ref17,ref18} remain the closest content-matched but Western-derived, Thai-unvalidated, family.

Layer 0 (a proposed extension)'s anonymous scored self-check has real precedent, though the precedent does not show benefit. The most rigorously tested example, I-DECIDE, is a masked two-arm RCT ($n=422$, 79--80\% 12-month retention) built from an explicit theoretical model and published protocol \cite{ref119,ref121}, with the methodological and ethical challenges of running the web-based trial documented separately \cite{ref122}; at 6 and 12 months it found no improvement in self-efficacy, with scores in fact slightly lower in the intervention arm than control (mean difference 1.3, 95\% CI 0.3--2.3 at 6 months; 1.6, 95\% CI 0.5--2.7 at 12 months), and no between-group difference in depressive symptoms \cite{ref120}. The myPlan app family, tested across protocol, longitudinal and sub-population studies \cite{ref123,ref126,ref127,ref128}, reported a significant reduction in reproductive coercion at 12 months ($p=.019$), but its own reported $p$-value for suicide-risk improvement ($p=.46$) does not support that outcome as significant despite wording in the source abstract implying otherwise \cite{ref124}; a bystander sub-study and teen adaptation extend the programme without changing this picture \cite{ref125,ref126}. A double-blind Canadian trial of a tailored online safety and health tool (planned $n=450$ in the protocol \cite{ref129}, $n=462$ enrolled in the completed trial \cite{ref130}) is the closest design analogue to a full self-check-plus-routing model, and a user-centred design report catalogues the wider suite of screening/safety apps this field has produced \cite{ref131}. The only Thailand-based example, a Thai adaptation of iCanPlan, has so far been evaluated by experts only (usability 4.93/5, cultural appropriateness 5.00/5), not by users \cite{ref132}; a military-partner survey and two EHR/computer-assisted screening trials extend the pattern without adding a premarital, scored, self-directed instrument \cite{ref133,ref134,ref135}. A Kenya-adapted risk instrument shows a scored tool can be validly tailored to a non-Western setting (AUC 0.78) \cite{ref136}, but a systematic review of ICT-based IPV interventions notes safety, equity and unintended-consequence research is ``virtually non-existent'' \cite{ref137}, and a review of web-based/mHealth IPV interventions finds computer-based/self-administered screening is the most common, not the best-proven, approach \cite{ref138}. This literature positions Layer 0 as feasible and acceptable, not validated: no record here is evidence a premarital scored self-check changes recognition, help-route knowledge or enacted choice.

\subsection{Multicultural clinical safety, language access, and Islamic medical ethics}
This is the evidence most directly relevant to \tphe{}'s cultural-intelligibility safeguard, and it is consistently cautionary. Cultural-competence training for health professionals repeatedly shows provider-level knowledge or attitude gains without a corresponding patient-outcome gain: a Cochrane review of 5 RCTs (337 professionals, 8{,}400 patients) rated the evidence low-quality, with some improvement in involvement/understanding but no reliable clinical-outcome effect, and adverse outcomes not assessed in any trial \cite{ref55}; a narrower review (7{,}879 records to 5 outcome-reporting studies) found no significant patient-outcome improvement in any \cite{ref56}; and Indigenous cultural-safety training across four high-income countries found outcomes overwhelmingly provider-level, with weak community-involvement and clinical-outcome evidence \cite{ref57}. Cultural humility \cite{ref58} and cultural safety --- a shift ``from provider control to patient agency'' \cite{ref59} --- are the constructs this paper adopts for the cultural-intelligibility safeguard, given the pattern above in which content-delivery training alone has not been shown to close this outcome gap. Language-access evidence sharpens the case for interpretation as distinct from cultural content: language discordance is associated with worse outcomes in over half of studies reviewed (84{,}750 participants) \cite{ref60} and with higher unplanned readmission/ED revisit \cite{ref61}; professional interpretation improves patient satisfaction and communication quality relative to ad hoc family/friend interpretation \cite{ref62}, and produces fewer linguistic errors than ad hoc interpretation \cite{ref63}; and a Canadian cohort ($n=124{,}583$) found language-concordant care associated with 36\% lower major cardiovascular events among allophone speakers \cite{ref64}. Thailand-specific evidence situates this in \tphe{}'s delivery context: interpretation/cultural-mediation for migrant health in Thailand faces budget, training and system-status constraints \cite{ref65}; migrant-health coverage in Thailand is under-studied in reproductive/maternal domains relative to general health \cite{ref66}; a Thai maternity-care ethnographic study named ``the system is in control,'' with insensitivity to cultural practice a subtheme \cite{ref67}; Thai Muslim NCD patients in the rural south report medication-adherence barriers (including Ramadan-related non-adherence) tied to social/cultural context \cite{ref68}; and Muslim mothers describe faith, patience and family/professional support as their coping structure after their children's cardiac surgery \cite{ref69}. Migrant/Muslim-women SRH-access reviews consistently name information, language and cultural-difference barriers \cite{ref70,ref20}, while quasi-experimental tailored delivery in southern Thailand nearly doubled cervical-screening uptake \cite{ref71} --- a single behavioural endpoint, not evidence that tailoring transfers to \tphe{}'s decision-control outcome.

Tier-one Islamic bioethics literature gives \tphe{}'s ethical frame real companions in the medical literature, without becoming evidence the frame, or \tphe{}, works. Islamic ethico-legal reasoning has been mapped onto the four principles of biomedical ethics \cite{ref225} and surveyed as a field spanning virtue ethics to jurisprudential rulings \cite{ref228,ref230}; empirical work with Muslim clinicians describes real moral distress reconciling Islamic and secular institutional ethics and calls for religious-competence training and peer support \cite{ref227}, while cross-gender-interaction norms are addressed separately \cite{ref226}, and end-of-life analyses treat consent as a shared, family-involved process that does not displace individual assent \cite{ref229}. Two Journal of IMA records, admitted under a labelled professional-association carve-out rather than as leading-journal evidence, frame informed consent as properly incorporating religious/cultural values \cite{ref241} and frame Qur'an-guided physician character as trustworthy, accountable professional conduct \cite{ref242}. Faith-integrated delivery can carry a measurable clinical effect: a cluster-randomised mosque-based diabetes-prevention trial reduced 12-month type 2 diabetes incidence from 17.1\% to 9.8\% (adjusted hazard ratio 0.75, 95\% CI 0.60--0.95, $P=.02$) \cite{ref231}, and a mosque-based, lay-led Islamic Trauma Healing programme for refugee PTSD showed effect sizes of $d=-0.67$ (PTSD severity), $d=-0.66$ (depression) and $d=0.71$ (well-being), maintained at 12-week follow-up \cite{ref232} --- though two published errata (\emph{JAMA Netw Open} 2024;7(9):e2440198 and 2024;7(10):e2446257) mean the corrected article, not the original abstract alone, should be checked before any number from it is repeated. A religiosity-anxiety meta-analysis found a consistent protective association (mean $r=-0.22$ across 10 studies) \cite{ref238} (further tier-one companions on vaccination decision-making, critical-illness family needs and faith-based infectious-disease framing: \cite{ref233,ref234,ref243}). This literature also complicates any simply-protective reading: a systematic review of gender-based violence among migrant Muslim women found the husband's control over the woman's life a risk factor for violence at the microsystem level, even where spirituality was protective at the individual level \cite{ref235}, and a systematic review of self-immolation found marital and familial conflict a principal driver in some Muslim-majority settings \cite{ref239} (further tier-one companions on sexual/reproductive health, assisted reproduction and religion-and-health: \cite{ref236,ref237,ref240}). No indexed FIMA or IMANA premarital guideline and no Thai Muslim premarital trial exist in this record. This is what the abstract's ethical terms are declared alongside, in their own field's own words, with a plain gloss attached --- the frame's companion literature, never evidence that the frame prevents harm.

\subsection{Thailand and the southern border provinces}
Direct Thai evidence remains sparse and geographically uneven. A four-region national survey ($n=2{,}462$ married/cohabiting women) found roughly 15\% lifetime psychological/physical/sexual IPV and controlling behaviours ranging 4.6--29.3\% \cite{ref72}; clinic-based studies found 21\% past-year IPV across two Bangkok gynaecology clinics (IPV associated with STI, aOR 2.65) \cite{ref73} and 13.1\% ever-abused among 475 pregnant Thai women \cite{ref74}. Provider-side evidence --- from \emph{northeast} Thailand, not the southern border --- found providers valued IPV screening but cited cultural-norm, training and service-condition barriers \cite{ref75}, a caution against generalizing provider attitudes across Thai regions. Southern-border-specific evidence is more extensive: among 700 Muslim adolescents on the southern border, 9\% were sexually experienced, with condom non-use, STI and repeat pregnancy concentrated in that subgroup \cite{ref76}; among 360 Muslim army conscripts in the three deep southern provinces, 78.7\% showed poor HIV/STI-transmission knowledge, examined further through an Islamic-values framing of sexual knowledge in a related conscript sample \cite{ref77,ref78}. Mental-health evidence from the conflict-affected southernmost provinces is consistent in direction across cohorts: migration and economic stress were associated with more psychiatric symptoms and religiosity with fewer among 2{,}053 Muslim adults \cite{ref79}; among youth 18--24, parental absence and travel-safety worry were associated with more psychiatric symptoms, with religious practice protective for men \cite{ref80}; among 387 adolescents, 7.2\% screened at risk of depression, with family dysfunction and PTSD strongly associated with depression risk \cite{ref81}; and young adults wishing to migrate for work reported more psychiatric symptoms and lower happiness, with chronic violence linked to education, employment and settlement plans \cite{ref82}. The same religious-interpretation and trust structure recurs in vaccine-acceptance evidence: per interviewed health workers, perceived religious prohibition, spousal disagreement, fear of side effects and religious-leader views were the leading barriers to MMR-vaccine acceptance among Muslim parents \cite{ref83}, and childhood-immunization beliefs were shaped by self-interpretation, misinformation and halal concerns, tempered by reliable information and community health-network leaders \cite{ref84}. No study identified in either round evaluates a premarital-education programme, of any design, with a Thai Muslim or Thai intercultural population.

\subsection{Culture as a health determinant, humility, and self-regulation}
Framed at the widest level, culture-and-health scholarship supplies theoretical warrant, not efficacy evidence, for treating culture as a determinant \tphe{}'s safeguards must hold rather than resolve. The Lancet Commission on Culture and Health argues culture shapes the meaning of health, behaviour, trust and communities of care \cite{ref85}; a companion Lancet series maps mechanisms --- separation and hierarchical power \cite{ref86}, and pathways through health systems and communities \cite{ref87} --- by which racism and discrimination affect health, and a 2026 follow-up renews calls for inclusive services, equitable access, rights and dignity for people on the move \cite{ref88}. Evidence on whether \emph{adapting} an intervention to a specific culture adds benefit beyond a generically delivered one is more equivocal than commonly assumed: a meta-analysis of 67 RCTs of culturally adapted psychological interventions for people of Chinese descent found a medium overall effect but \emph{no significant difference} between culturally modified and culturally specific adaptation types \cite{ref89} --- evidence against assuming any particular adaptation strategy is superior, not against cultural adaptation itself. Two structured-social-support RCTs outside the marriage-transition setting offer an analogy, not evidence, for the mechanism \tphe{}'s safeguards are built around: a Nepal social-capital-building programme was associated with lower incident major depressive disorder (7.4\% vs 12.9\%) \cite{ref90}, and a UK family-wellbeing RCT (674 participants, 62\% ethnic-minority) improved parental mental wellbeing post-intervention and at 6 months \cite{ref91}. The pattern is read here as suggestive that structured social support, not content delivery alone, may be the active ingredient --- a hypothesis adopted by analogy, not one either trial tested directly.

The constructs the model uses to gloss \emph{faqr} and \emph{taqwa} --- humility, compassion and self-regulation --- are real, measured psychological variables, but the evidence sits mostly in general and clinical populations, not premarital or Muslim ones, and the sequence to safety is not self-executing. Intellectual humility predicts more constructive responses to conflict \cite{ref203}, greater empathic accuracy toward outgroup members \cite{ref204}, and greater adherence to expert recommendations even controlling for political affiliation \cite{ref205}; cultural humility has a developed theoretical basis \cite{ref206,ref207} and institutional training improves attitudes and knowledge, with the authors themselves cautioning that a single training session is insufficient \cite{ref208}. Self-compassion has a large protective association with psychopathology ($r=-0.54$, 95\% CI $-0.57$ to $-0.51$) \cite{ref209}, and compassion-focused therapy improves self-compassion and reduces symptoms in trial evidence \cite{ref210}, but its relational benefit is conditional: among men, self-compassion predicted constructive repair only when conscientiousness was high, and the association reversed at low conscientiousness \cite{ref211}; two further studies link self-compassion to relationship satisfaction and conflict resolution in general and Black heterosexual-couple samples \cite{ref212,ref213}. Empathy shows the same double edge: working memory was associated with less physical abuse perpetration, an effect mediated by affective empathy \cite{ref214}, but in one clinical sample of men, emotional empathy was positively associated with perpetrated psychological IPV \cite{ref215} --- restraint in the use of power, not empathy alone, is the load-bearing step. Self-regulation deficits (inattention, impulsivity) are cohort-study risk factors for adult IPV perpetration, and existing treatment programmes are reported to have minimal impact to date \cite{ref216}; a systematic review traces the broader psychosocial pathway from childhood intra-familial victimisation to adult relationship violence \cite{ref217}. Religiosity and spirituality were protective for general interpersonal aggression but \emph{not} specifically for domestic violence in a 67-study systematic review and meta-analysis \cite{ref220}, complicating any simple ``faith protects'' reading (further tier-one companions extending without resolving this picture: \cite{ref221,ref223}). Trust and dyadic-maintenance scales \cite{ref218,ref219}, a humility-forgiveness mediation study \cite{ref202}, and the only premarital analogues in this literature --- an 8-couple feasibility pilot of a culturally adapted couple-learning programme \cite{ref100} and a premarital HIV-counselling cluster RCT \cite{ref224} --- round out the picture. Humility and compassion are evidence-bearing constructs by analogy from therapy and general populations; they position R1's ethical language, they do not supply evidence that R1 works in this setting.

\begin{table*}[t]
\centering\footnotesize
\caption{Evidence map for propositions P1--P7 (reference numbers as in the bibliography; each entry carries its PMID).}
\label{tab:evidencemap}
\begin{tabularx}{\textwidth}{@{}p{0.9cm}YYY@{}}
\toprule
& \textbf{Relevant / complicating records} & \textbf{Evidence level} & \textbf{What remains untested}\\
\midrule
P1 & \emph{Relevant:} burden/instrument base \cite{ref1,ref2,ref5,ref16,ref18,ref92}; norms (not knowledge) drove 70--95\% of the SASA! effect \cite{ref26,ref25}; decision-aid/shared-decision-making analogue \cite{ref14,ref15}. \emph{Complicating:} RCT-only aggression effect not convincing \cite{ref6,ref33,ref7,ref8}; uneven knowledge/literacy-to-outcome transfer \cite{ref34,ref35}; Muslim identity does not uniformly predict lower autonomy \cite{ref93}. & Mixed: burden data and single-time-point instrument-validity studies are strong; the causal claim that a governance layer \emph{adds} decision control beyond content, and that decision-quality measures capture it, is untested for any premarital-education format. & Whether a governance/safeguard layer added to premarital content produces a measurable decision-control effect distinct from content alone, using any instrument (existing or adapted), has not been tested.\\
\addlinespace
P2 & \emph{Relevant (adjacent precedent only --- no cited record tests a private channel specifically inside premarital/marriage education):} \cite{ref94,ref95,ref96,ref97,ref98,ref99}. \emph{Complicating:} even validated instruments substantially undercount coercive control \cite{ref51}. & Mixed: instrument-validation and single trials; retention/attrition data are a direct caution. & Whether a private channel embedded specifically in premarital education (not a clinical encounter) reliably surfaces concealed coercion has not been tested.\\
\addlinespace
P3 & \emph{Relevant:} \cite{ref100,ref101,ref27}; tailored edutainment nearly doubled cervical-screening uptake, a single behavioural endpoint \cite{ref71}. \emph{Complicating:} \cite{ref102,ref103,ref19}; cultural-competence training repeatedly improves provider knowledge/attitude without patient-outcome gain \cite{ref56,ref57,ref55}; no efficacy difference between adaptation types in one large meta-analysis \cite{ref89}. & Mixed to low: mostly single small trials, pilot studies, and systematic reviews of provider-level rather than patient-level outcomes; no controlled comparison of tailored vs untailored delivery on a safety outcome. & Whether cultural tailoring can be operationalized so that it improves engagement without ever being permitted to override a consent/safety threshold has not been tested in any premarital-education trial.\\
\addlinespace
P4\slash STOP & \emph{Relevant/analogous (not effectiveness support):} \cite{ref45,ref46,ref47,ref48,ref9,ref10,ref43}. \emph{Complicating:} \cite{ref49,ref50}. & Bounded: the strongest single review here is a meta-analysis of only 6 small studies. & No study has tested \tphe{}'s specific stop rule (R4) against a graduated-response alternative; the assessment mechanism that would trigger it is not specified or validated in \tphe{}'s own materials.\\
\addlinespace
P5 & \emph{Relevant:} 2025 USPSTF evidence report and recommendation statement \cite{ref38,ref39,ref40}; referral pathway \cite{ref42,ref43,ref44,ref104,ref105,ref106,ref107}; language-concordant referral \cite{ref60,ref62,ref63,ref61,ref64}. \emph{Complicating:} provider training improves knowledge/readiness, not reliably behaviour \cite{ref108,ref109,ref41}. & Strong for the general principle; weak/absent for \tphe{}'s own referral pathway, and untested for whether it meets language/cultural-accessibility standards in a Thai setting. & Whether \tphe{}'s proposed referral pathway meets timeliness/accessibility/follow-through \emph{and} language/cultural-accessibility thresholds in a Thai setting has not been tested.\\
\addlinespace
P6 & \emph{Relevant:} \cite{ref110,ref28,ref111,ref112}; a faith-leader trial found no effect on its primary outcome (secondary IPV effects only) \cite{ref29}; adjacent-field confirmation: screening uptake ``suboptimal'' despite an established programme \cite{ref12}; knowledge gains consistent, uptake inconsistent \cite{ref13}; RHL gains uneven across dimensions \cite{ref34}; knowledge/attitudes rose, no clinical outcomes reported \cite{ref35}; provider training raises knowledge/attitude without patient-outcome gain \cite{ref55,ref56,ref57}. & Strong: multiple RCTs and systematic reviews across two literatures now directly demonstrate this knowledge/behaviour and process/outcome dissociation. & Whether this dissociation pattern specifically occurs in \tphe{}'s own programme has not been tested --- it is now well-grounded across two adjacent literatures, not yet observed in \tphe{}.\\
\addlinespace
P7 & \emph{Relevant:} \cite{ref113,ref114,ref31,ref115,ref116}; methodological scaffolding for future revision cycles: MRC complex-intervention framework \cite{ref52}, ADAPT guidance \cite{ref53}, MRC process-evaluation guidance \cite{ref54}. & Weak to moderate: mostly programme-description and narrative-review evidence, plus general-purpose methodological guidance, not outcome trials of repair mechanisms themselves. & Whether \tphe{} has, or will build, an actual mechanism (audit cycle, feedback loop, revision trigger) for repairability has not been specified or tested.\\
\bottomrule
\end{tabularx}
\end{table*}

\subsection{Gap statement}
\tphe{} is offered as a programme theory --- a structured account of the mechanism by which a governance layer of five rules (R1--R5), including a stop--assess--refer--record rule (R4), is hypothesized to preserve or restore each participant's decision control at the marriage transition --- and not as a validated or effective intervention. No component of \tphe{} has been tested, in any population, against the outcomes its own failure criterion names as necessary (informed choice, coercion recognition, reproductive and mental-health literacy, safe help-seeking, accountable referral). This review adds two findings that sharpen, without resolving, the case for testing: first, that a premarital contact point can plausibly be built and delivered at population scale (drawing on preconception medicine and premarital genetic screening as the nearest analogue, not a tested precedent for \tphe{} itself), while knowledge gained at that contact point does not reliably convert into the behavioural or decision outcome desired, in either the screening or the relationship-education literatures; second, that cultural-competence and cultural-tailoring interventions elsewhere in health care repeatedly raise provider- or participant-level knowledge and attitudes without a demonstrated patient-outcome gain, and that no adaptation strategy has been shown superior to another in the one large trial that tested this directly. No existing instrument for any of \tphe{}'s target outcomes has been validated in a Thai or Thai Muslim population, and the decision quality/shared-decision-making measurement family this paper proposes as the nearest analogue for ``decision control'' is itself an unvalidated adaptation, not a borrowed instrument. \tphe{} should be read, and should be tested, as a falsifiable synthesis awaiting empirical evaluation --- never cited, by its authors or by others, as evidence that premarital education of this design improves, prevents, or reduces any outcome.
